\documentclass[10pt,a4paper]{amsart}
\usepackage[margin=19mm]{geometry}
\usepackage{amsmath,amssymb,mathtools}
\usepackage[scr=rsfs,cal=euler]{mathalfa}
\usepackage{xfrac}
\usepackage[unicode,hidelinks,bookmarks=false]{hyperref}
\newcommand{\ie}{i.e.\ }
\newcommand{\cf}{cf.\ }
\newcommand{\resp}{respectively\ }
\newcommand{\Subsection}{\subsection}
\providecommand{\nobreakdash}{\nobreak\hbox{-}}
\setlength{\emergencystretch}{2em}
\multlinegap0pt
\usepackage{bm}
\usepackage[normalem]{ulem}
\usepackage{amssymb}
\usepackage{xcolor}
\usepackage{booktabs}
\usepackage{algorithm}\usepackage{algpseudocode}
\usepackage{mathtools}
\usepackage{dsfont}
\usepackage[shortlabels]{enumitem}
\usepackage{bbm}
\usepackage{tikz}
\usepackage{cancel}
\usepackage{float}
\usepackage{tcolorbox}
\newtheorem{thm}{Theorem}
\title{{\Large \bf Modified Patankar Semi-Lagrangian Scheme for the Optimal Control of Production-Destruction systems}}
\author{Simone Cacace, Alessio Oliviero, Mario Pezzella}
\date{Source: arXiv:2501.13085v1 (CC BY 4.0)}
\begin{document}
\maketitle

\noindent\emph{Source and licence.} {\Large \bf Modified Patankar Semi-Lagrangian Scheme for the Optimal Control of Production-Destruction systems}. Version arXiv:2501.13085v1 (CC BY 4.0), distributed by arXiv under the Creative Commons Attribution 4.0 International licence, http://creativecommons.org/licenses/by/4.0/, as recorded in the arXiv licence metadata for this exact version and shown on the abstract page https://arxiv.org/abs/2501.13085v1. Full licence notice: \texttt{licenses/arxiv-2501.13085-CC-BY-4.0.txt}.\par\smallskip

\noindent\textit{Editorial scope.} Counted unit: the article's thm solution (source0.tex lines 242--244). The article's complete original proof of it is delivered with it (source0.tex lines 245--262, one \texttt{proof} environment of 18 lines). No mathematics is written, repaired or re-derived: the statement and the proof are the article's own lines, in the article's own notation, and no retained line is substituted or re-flowed.  Retained uncounted as the source-local context the proof needs: lines 214--221.  Nothing was omitted from any retained span, so the package carries no omission record.  The retained lines cite 1 works, whose entries are transcribed verbatim from the article's own bibliography source in the e-print and delivered with short editorial labels recorded in the item's reference mapping.  No retained line contains a figure environment, an image-inclusion command or drawing code, so no image is delivered and the package declares no asset dependency.  Editorial formatting only, in addition: an AMS \texttt{amsart} preamble that restates the article's own declarations and macros used by the retained lines, namely the environments \texttt{thm} on separate counters, together with the standard packages those lines need.  The retained environments print in this document as Theorem 1 in order of appearance, whereas the article prints the same items with the numbering of its own full document.  The complete native source of the article as distributed in the arXiv e-print for this exact version is bundled unchanged as \texttt{sources/171991/source0.tex}.
\begin{equation}\label{eq:CPDS}
\left\{
\begin{aligned}
    \bm{y}'(t)&=\big(P(\bm{y}(t))\odot \mathcal{P}(\bm{\alpha}(t))-D(\bm{y}(t))\odot \mathcal{D}(\bm{\alpha}(t)) \big) \bm{e}, \\
    \bm{y}(t_0)&=\bm{y}^0, \qquad \qquad \qquad  \qquad \qquad \qquad  \qquad \qquad \qquad \qquad  t\geq t_0\geq 0,
\end{aligned}
\right.
\end{equation}

\begin{thm}[\textbf{Existence of the solution}]\label{thm:Existence of a solution}
    Assume that $P$ and $D$ are Lipschitz continuous functions on $\mathbb{R}^N$ and that $\mathcal{P},\mathcal{D}\in C^0(A).$  Then, for each $\bm{\alpha}\in\mathcal{A},$ the controlled production-destruction system \eqref{eq:CPDS} admits a unique solution. 
\end{thm}

\begin{proof}
    The CPDS \eqref{eq:CPDS} can be rewritten as $\bm{y}'(t)= \bm{F}(\bm{y}(t),\bm{\alpha}(t))$, with $\bm{y}(0)=\bm{y}^0,$  $\bm{\alpha}\in \mathcal{A}$ and
    \begin{equation}\label{eq:F}
        \bm{F} : (\bm{\mu},\bm{a})\in \mathbb{R}^N\times A \longrightarrow \big(P(\bm{\mu})\odot \mathcal{P}(\bm{a})-D(\bm{\mu})\odot \mathcal{D}(\bm{a}) \big) \bm{e}  \in \mathbb{R}^N.
    \end{equation}
    Let $L_P,$ $L_D$ be the Lipschitz constants of $P$ and $D$, respectively. Define 
    \begin{equation*}
        \bar{p}=\!\!\max_{\substack{\bm{a}\in A \\ 1\leq i,j\leq N}}|\mathcal{P}_{ij}(\bm{a})| \qquad \qquad \text{and} \qquad \qquad  \bar{d}=\!\!\max_{\substack{\bm{a}\in A \\ 1\leq i,j\leq N}}|\mathcal{D}_{ij}(\bm{a})|.
    \end{equation*}
    The hypotheses on the known functions imply that $\bm{F}\in C^0(\mathbb{R}^N\times A)$ and that, for each $\bm{\mu},\bm{\nu} \in \mathbb{R}^N$ and $\bm{a}\in A,$ 
    \begin{equation*}
        \begin{split}
            \|\bm{F}(\bm{\mu},\bm{a})-\bm{F}(\bm{\nu},\bm{a}) \| \leq \; & \hat{e}\|\big (P(\bm{\mu})-P(\bm{\nu}))\odot \mathcal{P}(\bm{a})  \|+\hat{e}\|\big (D(\bm{\mu})-D(\bm{\nu}))\odot \mathcal{D}(\bm{a}) \| \\
            \leq \; &\hat{e}k_N\left(\bar{p}L_P  +\bar{d}L_D\right)\|\bm{\mu}-\bm{\nu}\|,
        \end{split}
    \end{equation*}
    where $\hat{e}=\|\bm{e}\|$ and $k_N>0$ is a constant depending on $N.$ As a result, the function $\bm{F}$ in \eqref{eq:F} is Lipschitz continuous as well and all the hypotheses of the Carathéodory theorem (see \cite[Theorem 8.1]{FalconeFerretti}) are fulfilled. The existence of a unique solution $\bm{y}(t)$ to \eqref{eq:CPDS} is then assured.
\end{proof}

\begin{thebibliography}{99}
\bibitem[Falcon]{FalconeFerretti} M.~Falcone and R.~Ferretti. \newblock {\em Semi-Lagrangian Approximation Schemes for Linear and {Hamilton}--{Jacobi} Equations}. \newblock Society for Industrial and Applied Mathematics, 2013. \newblock \href {https://doi.org/10.1137/1.9781611973051} {\path{doi:10.1137/1.9781611973051}}.
\end{thebibliography}
\end{document}
